

\chapter{Diskussion}
\label{chap:diskussion}

\section{Nutzenbewertung}

Suchzeit
Wiederverwendung

\section{Wirtschaftliche Betrachtung}

Verbindung zu Kapitel 1.

\section{Grenzen}

\chapter{Grenzen und Ausblick}
\label{chap:grenzen_und_ausblick}

\section{Grenzen des aktuellen Systems}

Die Ergebnisse der Evaluierung zeigen, dass die größte Einschränkung des aktuellen Systems in der Variantenauflösung technisch ähnlicher Makrofamilien liegt.
Dies lässt sich auf die Struktur sowie die Eigenschaften der verfügbaren Datenbasis zurückführen.\parencite{LTR_Configuration_Test_Series}

Tabelle \ref{tab:grenzen_uebersicht} fasst die zentralen Grenzen des Systems und die daraus abgeleiteten Ansatzpunkte zusammen.

\begin{table}[H]
	\centering
	\setlength{\tabcolsep}{4pt}
	\renewcommand{\arraystretch}{1.1}
	\caption{Übersicht der zentralen Grenzen und Ansatzpunkte}
	\label{tab:grenzen_uebersicht}
	\begin{tabularx}{\textwidth}{>{\raggedright\arraybackslash}p{3.5cm} >{\raggedright\arraybackslash}p{4.0cm} >{\raggedright\arraybackslash}X}
		\toprule
		\textbf{Grenze} & \textbf{Beobachteter Befund} & \textbf{Ansatzpunkt} \\
		\midrule
		Geringer Support & 302 Kandidatenprofile basieren nur auf einem Trainingsartikel & Erweiterung der Datenbasis für seltene Ziel-Makros \\
		Variantenauflösung & Family Recall@1 liegt deutlich über Target Recall@1 & Spezialisierte Modellierung innerhalb der Makrofamilien \\
		Support-Bias & Häufige Ziel-Makros besitzen stabilere Kandidatenprofile & Ausgleich seltener Profile durch zusätzliche Beispiele und Profilpflege \\
		\bottomrule
	\end{tabularx}
\end{table}

\subsection{Geringer Support von Ziel-Makros}

Eine zentrale Einschränkung ergibt sich aus der Support-Verteilung der Kandidatenprofile. Von insgesamt 375 vorhandenen Kandidatenprofilen basieren 302 Profile auf nur einem Trainingsartikel.
Der Median-Support liegt bei 1 und der Mittelwert der Artikel pro Ziel-Makro \parencite{EPLAN-Makros} beträgt 2,62.

Ein einzelner Artikel ist dabei nicht in der Lage, einen ausreichenden Support für die technische Vielfalt innerhalb einer Makrofamilie zu bieten.
Die daraus entstandenen Kandidatenprofile sind somit von nur einem Artikel abhängig.

\subsection{Variantenauflösung innerhalb einer Makrofamilie}

Die Resultate der Evaluierung zeigen, dass das Modell die korrekte Makrofamilie deutlich zuverlässiger erkennt als das exakte Ziel-Makro \parencite{EPLAN-Makros}.
A08 als beste Konfiguration erzielt dabei einen Family Recall@1 von 87,35\,\%, wohingegen der Target Recall@1 lediglich 34,78\,\% beträgt.

Zudem zeigt die Fehleranalyse, dass in 133 Testfällen die korrekte Makrofamilie erkannt wird, jedoch nicht das exakte Ziel-Makro \parencite{EPLAN-Makros}.
Die größte Herausforderung liegt somit in der fehlerfreien Unterscheidung zwischen sehr ähnlichen Varianten innerhalb derselben Makrofamilie \parencite{LTR_Configuration_Test_Series}.

\subsection{Support-Bias häufiger Ziel-Makros}

Der Support-Bias bezeichnet einen Rankingvorteil von Ziel-Makros mit hohem Kandidaten-Support. Da die Kandidatenprofile dieser Ziel-Makros auf mehreren Trainingsartikeln beruhen, können dadurch breitere technische Merkmale abgedeckt werden und dadurch häufiger durch das Modell gefunden werden.

Die durchgeführten Fehleranalysen zeigen einen deutlichen Support-Bias. In 55,6\,\% der Fälle besitzt das fälschlich ausgewählte Ziel-Makro \parencite{EPLAN-Makros} einen höheren Support als das korrekte Ziel-Makro \parencite{EPLAN-Makros}.
Im Durchschnitt verfügt das falsch gewählte Ziel-Makro dabei über etwa 10,6 Artikel mehr als das korrekte Ziel-Makro \parencite{LTR_Configuration_Test_Series}.

Häufige Ziel-Makros \parencite{EPLAN-Makros} erzeugen dadurch stabilere Kandidatenprofile und bekommen somit einen strukturellen Vorteil gegenüber selteneren Varianten.



\section{Kritische Reflexion}


% chapter{Diskussion der Ergebnisse}\label{chap:diskussion_der_ergebnisse}

% \section{Ergebnisse des Validierungssplits}

% Die Ergebnisse des Validierungssplits werden im späteren Verlauf für die Auswahl der besten Modellkonfiguration verwendet.

% Tabelle \ref{tab:validation_ergebnisse} fasst die wichtigsten Rankingmetriken des Referenzlaufs A08 auf dem Validierungssplit zusammen. 

% \begin{table}[H] 
%   \centering 
%   \caption{Ergebnisse des Referenzlaufs A08 auf dem Validierungssplit} 
%   \label{tab:validation_ergebnisse} 
%   \begin{tabular}{lr} \toprule \textbf{Metrik} & \textbf{Wert} \\ 
%     \midrule 
%     Target Recall@1 & 0,3137 \\
%     Target Recall@3 & 0,4472 \\
%     Target Recall@5 & 0,5031 \\
%     Target Recall@10 & 0,6211 \\
%     Family Recall@1 & 0,8199 \\
%     Family Recall@10 & 0,8913 \\
%     MRR & 0,4121 \\
%     Hierarchical nDCG@10 & 0,6766 \\
%     \bottomrule 
%   \end{tabular} 
% \end{table}

% Die Ergebnisse zeigen, dass das Modell bereits auf dem Validierungssplit die korrekte Makrofamilie besser identifiziert als das exakte Ziel-Makro \parencite{EPLAN-Makros}.
% Während das korrekte Ziel-Makro \parencite{EPLAN-Makros} in 31,37\,\% der Anfragen auf Platz 1 steht, wird die korrekte Makrofamilie in 81,99\,\% der Anfragen auf die erste Position gesetzt.

% Durch den Target Recall@10 lässt sich erkennen, dass das Modell in 62,11\,\% der Anfragen das korrekte Ziel-Makro \parencite{EPLAN-Makros} innerhalb der Top-10-Vorschläge platziert.
% Interessant ist der Abstand zwischen Family Recall@10 (89,13\,\%) und Target Recall@10 (62,11\,\%). Dies deutet darauf hin, dass viele Fehlkategorisierungen innerhalb derselben Makrofamilie stattfinden.

% Der erzielte MRR-Wert von 0,4121 sowie das hierarchische nDCG@10 von 67,66\,\% zeigen zusätzlich, dass beim Validierungssplit das Modell in der Lage ist, eine sinnvolle Reihenfolge zu erzeugen, 
% dabei jedoch weiterhin Schwierigkeiten bei der korrekten Auflösung technischer Varianten besitzt, \parencite{LTR_Configuration_Test_Series}.

\section{Interpretation der Ergebnisse}

Nachdem die Modellauswahl abgeschlossen wurde, erfolgt die finale Evaluation des ausgewählten Modells auf dem unberührten Testsplit.

Tabelle \ref{tab:test_ergebnisse} zeigt die erzielten Ergebnisse des Referenzlaufs A08 auf dem Testsplit. 

\begin{table}[H] 
  \centering 
  \caption{Ergebnisse des Referenzlaufs A08 auf dem Testsplit} 
  \label{tab:test_ergebnisse} 
  \begin{tabular}{lr} \toprule \textbf{Metrik} & \textbf{Wert} \\ \midrule 
    Target Recall@1 & 0,3478 \\
    Target Recall@3 & 0,4783 \\
    Target Recall@5 & 0,5494 \\
    Target Recall@10 & 0,6798 \\
    Family Recall@1 & 0,8735 \\
    Family Recall@10 & 0,9368 \\
    MRR & 0,4534 \\
    Hierarchical nDCG@10 & 0,7449 \\
    \bottomrule 
  \end{tabular} 
\end{table}

Beim Testsplit zeigt sich ein ähnliches Bild wie beim Validierungssplit. Das Modell sagt bei 34,78\,\% der Anfragen das korrekte Ziel-Makro \parencite{EPLAN-Makros} auf Platz 1 voraus, während die korrekte Makrofamilie in 87,35\,\% der Anfragen auf die erste Position gesetzt wird.

Wie beim Validierungssplit ist auch hier der deutliche Abstand zwischen Family Recall@10 (93,68\,\%) und Target Recall@10 (67,98\,\%) zu erkennen. 

Der erreichte MRR-Wert von 45,34\,\% sowie die hierarchische nDCG@10 von 74,49\,\% bestätigen zusätzlich die Erkenntnisse aus dem Validierungssplit, \parencite{LTR_Configuration_Test_Series}.

\section{Ursachenanalyse der Rankingfehler}

Um die Modellleistung detailliert erfassen zu können, werden die Top-1-Ergebnisse des Testsplits analysiert. Dabei werden die fachlichen Beziehungen zum tatsächlichen Ziel-Makro untersucht.

Tabelle \ref{tab:fehleranalyse_top1} zeigt die Verteilung der einzelnen Fehlerkategorien. 

\begin{table}[H] 
  \centering 
  \caption{Top-1-Analyse der Vorhersagen auf dem Testdatensatz} 
  \label{tab:fehleranalyse_top1} 
  \begin{tabular}{lrr} 
    \toprule 
    \textbf{Kategorie} & \textbf{Anzahl} & \textbf{Anteil} \\ 
    \midrule 
    Exaktes Ziel-Makro \parencite{EPLAN-Makros} & 88 & 34,78\,\% \\
    Gleiche Familie, falsches Target & 133 & 52,57\,\% \\
    Gleiches Familienpräfix & 12 & 4,74\,\% \\
    Andere Familie & 20 & 7,91\,\% \\
    \bottomrule 
  \end{tabular} 
\end{table}

Die Fehleranalyse bestätigt die Beobachtungen aus den Metriken. Von insgesamt 253 Testanfragen werden 133 Fälle einer gleichen Familie, aber einer falschen Variante zugeordnet, was 52,57\,\% entspricht.
In 7,91\,\% der Fälle wird eine andere Familie vorhergesagt. Das zeigt, dass das Modell die grundlegenden technischen Zusammenhänge zwischen Artikeln und Makrofamilien zuverlässig erlernt hat.
Die Analyse zeigt, dass die zentrale Herausforderung die Identifikation des korrekten Ziel-Makros \parencite{EPLAN-Makros} innerhalb der richtigen Makrofamilie ist. Dafür muss die Variantenauflösung innerhalb der Makrofamilie weiter verbessert werden.
Die große Differenz zwischen dem Target Recall@1 und dem Family Recall@1 unterstreicht diese Herausforderung zusätzlich.
Dies wird im Kapitel \ref{chap:grenzen_und_ausblick} genauer analysiert, \parencite{LTR_Configuration_Test_Series}.

% \section{Rangstatistische Analyse}

% Eine weitere wichtige Analyse ist die statistische Auswertung der erzeugten Ranglisten.


% Die wichtigsten Rangstatistiken des Testdatensatzes werden in Tabelle \ref{tab:rangstatistik} dargestellt.

% Der Mean Rank ist der arithmetische Mittelwert der Rangpositionen des korrekten Ziel-Makros. Der Median Rank bezeichnet die Rangposition, unter beziehungsweise über der jeweils 50\,\% der Fälle liegen. Das $p$-Perzentil bezeichnet den Rang, den $p$ Prozent der Fälle nicht überschreiten.

% \begin{table}[H] 
%   \centering 
%   \caption{Rangstatistische Kennzahlen des Testsplits} 
%   \label{tab:rangstatistik} 
%   \begin{tabular}{lr} 
%     \toprule 
%     \textbf{Kennzahl} & \textbf{Wert} \\ 
%     \midrule 
%     Erwarteter Zielrang (Mean Rank) & 23,96 \\
%     Median Rank & 4 \\
%     90\,\%-Perzentil Rank & 61,6 \\
%     Maximaler Rang & 360 \\
%     \bottomrule 
%   \end{tabular} 
% \end{table}

% Der Median Rank von 4 zeigt, dass das korrekte Makro \parencite{EPLAN-Makros} in der Hälfte der Testanfragen unter den ersten vier Positionen der Rangliste erscheint.
% Der mittlere Rang (Mean Rank) von 23,96 zeigt gleichzeitig, dass es einzelne schwierige Anfragen gibt, die deutlich schlechter eingeordnet werden.
% Diese Annahme wird durch das 90\,\%-Perzentil von 61,6 bestätigt. Dabei sitzt das Makro \parencite{EPLAN-Makros} in 90\,\% der Testanfragen unter den ersten 61,6 Positionen der Rangliste.

% Auch die Rangstatistiken zeigen, dass das Modell in der Lage ist, die korrekte Makrofamilie in den meisten Fällen weit oben in der Rangliste zu platzieren, \parencite{LTR_Configuration_Test_Series}.

% \section{Generalisierungsfähigkeit}

% Die Generalisierungsfähigkeit bedeutet, dass das Modell in der Lage ist, zuverlässige Rankings für nicht im Training verwendete EPLAN-Artikel zu erzeugen.
% Zur Beurteilung der Generalisierungsfähigkeit des Modells wurde eine zusätzliche Leave-One-Out-Kontrollanalyse durchgeführt. Bei diesem Verfahren wird jeweils ein Beispiel zurückgehalten und anhand der verbleibenden Beispiele bewertet. Projektspezifisch wird für jeden Trainingsartikel das Kandidatenprofil seines eigenen Ziel-Makros ohne diesen Artikel neu gebildet und anschließend bewertet. Das trainierte Modell bleibt dabei unverändert; die Analyse dient ausschließlich diagnostischen Zwecken. Um eine Verfälschung der Ergebnisse der Trainingsanalyse zu verhindern, wird bei einer Anfrage genau dieser Artikel aus der Anfragemenge ausgeschlossen.


% Tabelle \ref{tab:loo_vergleich} stellt die Ergebnisse der Leave-One-Out-Analyse (LOO) sowie der Validierungs- und Testergebnisse gegenüber.

% \begin{table}[H] 
%   \centering 
%   \caption{Vergleich von regulärer Trainingsbewertung, Leave-One-Out-Analyse, Validation und Test} 
%   \label{tab:loo_vergleich} 
%   \begin{tabular}{lrrrr} 
%     \toprule 
%     \textbf{Metrik} & \textbf{Train (LOO)} & \textbf{Validation} & \textbf{Test} \\ 
%     \midrule 
%     Recall@1 & 0,1870 & 0,3137 & 0,3478 \\
%     Recall@3 & 0,3028 & 0,4472 & 0,4783 \\
%     Recall@10 & 0,4695 & 0,6211 & 0,6798 \\
%     MRR & 0,2864 & 0,4121 & 0,4534 \\
%     Family Recall@1 & 0,7246 & 0,8199 & 0,8735 \\
%     Hierarchical nDCG@10 & 0,5945 & 0,6766 & 0,7449 \\
%     Mean Rank & 25,93 & 23,96 & 23,96 \\
%     Median Rank & 13,0 & 4,0 & 4,0 \\
%     90\,\%-Perzentil & 62,0 & 61,6 & 61,6 \\
%         \bottomrule 
%     \end{tabular} 
%     \end{table}

% Es zeigt sich, dass die Leave-One-Out-Bewertung eine strengere Evaluationsmethode als die Bewertung durch den Validierungs- und Testdatensatz darstellt.
% Die Ursache dafür liegt im Support der Kandidatenprofile. Von den vorhandenen 375 Kandidatenprofilen basieren insgesamt 302 Profile ausschließlich auf einem Trainingsartikel.
% Wenn dieser Artikel aus der Anfragemenge ausgeschlossen wird, enthält das entsprechende Kandidatenprofil keine unterstützenden Trainingsdaten mehr und kann dadurch keinen Beitrag zur korrekten Vorhersage leisten.

% Bei den Validierungs- und Testdaten tritt dieser Effekt nicht auf, da durch die im Datensplit erzwungene Trainingsabdeckung für jedes Ziel-Makro \parencite{EPLAN-Makros} mindestens ein Artikel in jedem Split vorhanden ist, der nicht identisch zu den anderen beiden ist.


